% Discussão (S03-09, #128). Números vindos de analysis/ (ver resultados.tex).
\section{Discussão}
\label{sec:discussao}

Esta seção confronta as hipóteses da Seção~\ref{sec:introducao} com os 300
repositórios analisados. As hipóteses foram escritas antes de observar os
resultados. A Tabela~\ref{tab:hipoteses} resume o veredito, e os parágrafos
seguintes discutem o que ele significa e o que não permite concluir.

\begin{table}[ht]
\centering
\caption{Hipóteses prévias e resultados (300 repositórios).}
\label{tab:hipoteses}
\small
\begin{tabular}{lll}
\hline
Hipótese & Resultado & Veredito \\
\hline
H01 frequência heterogênea, abaixo de 1/semana & 89\% abaixo de 1/semana & Sustentada \\
H02 lead time (a) $>$ (b) & 250,8 h $>$ 169,4 h; 73\% dos repositórios & Sustentada \\
H03 CFR (a) e (b) divergem & medianas 0,054 e 0,292 & Sustentada \\
H04 recuperação assimétrica & mediana 4,9 h, Q3 21,5 h & Sustentada \\
H05 associação fraca ou não positiva & $\rho$ = 0,182 e 0,567 & \textbf{Refutada} \\
H06 diferença em ao menos uma métrica & 2 de 12 após Holm, ambas em lead time & Sustentada, só em lead time \\
H07 mudança de categoria nas variantes & 28,5\%--44,4\% nos pares comparados & Sustentada \\
\hline
\end{tabular}
\end{table}

\textbf{Contra o esperado.} A H05 previa associação fraca ou não positiva entre
frequência e CFR, e os dados mostram associação positiva com as duas variantes.
Isso contraria a leitura do DORA de que velocidade e estabilidade andam juntas.
Duas ressalvas impedem concluir que publicar mais gera mais falhas. A primeira
é que o CFR (b) tem a associação mais forte ($\rho = 0{,}567$) justamente por
construção: quem publica várias releases por semana tem mais chance de ter uma
\textit{patch} nos sete dias seguintes, que a heurística marca como corretiva.
Com o CFR (a), que não depende do ritmo de releases, a associação é fraca
($\rho = 0{,}182$), embora significativa com os 300 repositórios. A segunda é
que a amostra reúne projetos populares e muito ativos, em que frequência e CFR
podem crescer juntos por causa do tamanho do projeto e do volume de CI. É uma
associação entre postos, sem controle de confundidores.

\textbf{Efeito só no lead time.} Na RQ06, apenas 2 dos 12 testes são
significativos após Holm, e os dois envolvem o lead time (a): linguagem
($\varepsilon^2 = 0{,}099$) e quartil de contribuidores
($\varepsilon^2 = 0{,}053$). Frequência, CFR (a), recuperação e estrelas não
diferem. Projetos em TypeScript e Python têm lead time mediano menor do que os
em Rust e C++, e projetos com mais contribuidores entregam de forma mais
espaçada. Esses padrões são descritivos. Os testes não identificam quais pares
de grupos diferem, e linguagem e número de contribuidores podem refletir o tipo
de projeto, que não foi controlado. Alguns grupos de linguagem têm só 10 a 13
repositórios, e a ausência de significância nas demais combinações não
demonstra ausência de efeito.

\textbf{O CFR (b) está inflado.} Na validação de avaliador único, a heurística
tem recall de 0,870 e precisão de 0,238 (F1 = 0,374), abaixo do limiar de 0,70
do protocolo. Ela perde poucas releases corretivas (6 de 46), mas rotula como
corretivas muitas que não são (apenas 40 das 168 marcadas foram confirmadas).
Por isso a mediana do CFR (b), 0,292, deve ser lida como limite superior, não
como taxa de falha. A validação tem limites sérios: um só avaliador, 283
releases em 58 repositórios, e a rodada de consenso prevista não foi feita. Não
é possível separar o erro da heurística dos critérios pessoais do avaliador.

\textbf{Os proxies divergem.} Na classificação C1, nenhum repositório é Elite em
frequência e 136 (45\%) são Low, enquanto 230 (77\%) são Elite em CFR (a). A
categoria geral, a mediana das quatro pontuações, esconde essa oposição e se
concentra em Medium e High. A classificação depende da métrica de falha
escolhida, e a sensibilidade a essa escolha é o objeto da RQ07.

\textbf{Outliers.} Frequência, lead time e recuperação têm caudas longas
(lead time (a) com Q3 de 756 h, recuperação com Q3 de 21,5 h). Foram mantidos
em todas as análises, com mediana, IQR e testes por postos, e as figuras usam
escala \textit{symlog}. Testes por postos e medianas são pouco sensíveis a poucos
repositórios extremos, ao contrário da média, que seria distorcida.

\textbf{Limites de causalidade e de construto.} Release não é deploy, falha de
CI não é falha em produção, e a heurística de release corretiva é um proxy.
Nenhum resultado aqui mede incidentes reais. Todas as associações são
observacionais, e a amostra, de projetos populares de código aberto, não
representa equipes comerciais nem projetos menos populares.

\textbf{RQ07 e trabalho futuro.} A categoria geral muda em 44,4\% dos 259
repositórios comparáveis entre C1 e C2, em 28,5\% dos 260 entre C1 e C3, e em
41,0\% dos 261 entre C2 e C3. A concordância ponderada linear é,
respectivamente, $\kappa = 0{,}311$, $0{,}625$ e $0{,}428$. Esses resultados
mostram sensibilidade descritiva às variantes nesta amostra, não um efeito
causal da escolha da unidade de entrega. Em particular, a C2 usa o CFR (b),
cuja validação de avaliador único teve F1 = 0,374, abaixo do limiar de 0,70;
mudanças que envolvem C2 podem refletir erro da heurística. Permanecem como
trabalho futuro rotular a segunda rodada de validação com mais de um avaliador
e executar a replicação cruzada.
